\chapter{Scan Variables and Distributions}
\label{ch:variables}

\newthought{\yamlkey{Sampling.Variables}} lists the parameters you want to scan and tells \Jarvis\
how to draw values for each one. Every sampler reads the same variable format, so once you learn
it here it works everywhere.

\section{One variable at a time}
\label{sec:var-schema}

\yamlkey{Variables} is a list; each item defines \emph{one} variable:

\begin{yamlwin}
Sampling:
  Variables:
    - name: x1
      description: "A flat variable"
      distribution:
        type: Flat
        parameters:
          min: 0
          max: 10
\end{yamlwin}

\begin{description}[style=nextline, leftmargin=1.2em]
  \item[\yamlkey{name}] A short identifier, unique within the scan. You use this name later in
        likelihood and mapping expressions.
  \item[\yamlkey{description}] A one-line note on what the variable means. Optional but
        recommended---it shows up in diagnostics.
  \item[\yamlkey{distribution.type}] Which distribution to draw from (the next section lists
        them).
  \item[\yamlkey{distribution.parameters}] The numbers that distribution needs (for example
        \yamlkey{min} and \yamlkey{max}).
\end{description}

\section{The five distributions}
\label{sec:distributions}

\Jarvis\ supports the five one-dimensional distributions in Table~\ref{tab:distributions}. Each
has its own required parameters.

\begin{table}[ht]
  \footnotesize
  \begin{tabular}{@{}lll@{}}
    \toprule
    \textbf{\yamlkey{type}} & \textbf{Parameters} & \textbf{Use it when\ldots} \\
    \midrule
    \yamlkey{Flat}       & \yamlkey{min}, \yamlkey{max}        & values are equally likely across a range \\
    \yamlkey{Log}        & \yamlkey{min}, \yamlkey{max}        & the range spans orders of magnitude (\yamlkey{min}$>0$) \\
    \yamlkey{Normal}     & \yamlkey{mean}, \yamlkey{stddev}    & values cluster around a central value \\
    \yamlkey{Log-Normal} & \yamlkey{mean}, \yamlkey{stddev}    & the \emph{logarithm} clusters around a value \\
    \yamlkey{Logit}      & \yamlkey{location}, \yamlkey{scale} & a bell-like shape with heavier tails \\
    \bottomrule
  \end{tabular}
  \caption{The supported variable distributions.}
  \label{tab:distributions}
\end{table}

\paragraph{Flat (uniform)} Every value between \yamlkey{min} and \yamlkey{max} is equally likely:
$x \sim \mathcal{U}(\mathrm{min}, \mathrm{max})$. This is the default choice when you have no
prior preference.

\paragraph{Log (log-uniform)} Equally likely \emph{per decade}, so a scan covers
$10^{-3}$ as densely as $10^{3}$: $\ln x \sim \mathcal{U}(\ln\mathrm{min}, \ln\mathrm{max})$.
Requires \yamlkey{min}$>0$.

\paragraph{Normal (Gaussian)} A bell curve centred on \yamlkey{mean} with width
\yamlkey{stddev}: $x \sim \mathcal{N}(\mathrm{mean}, \mathrm{stddev})$. About 68\% of values land
within one \yamlkey{stddev} of the \yamlkey{mean}.

\paragraph{Log-Normal} A Normal distribution in $\ln x$ rather than $x$, useful for positive
quantities that vary multiplicatively: $\ln x \sim \mathcal{N}(\mathrm{mean}, \mathrm{stddev})$.

\paragraph{Logit (logistic)} A bell-like distribution with heavier tails than the Normal, set by
a centre \yamlkey{location} and width \yamlkey{scale}:
$x = \mathrm{location} + \mathrm{scale}\cdot\ln\!\big(\tfrac{u}{1-u}\big)$ for a uniform
$u \in (0,1)$, i.e. $x \sim \mathrm{Logistic}(\mathrm{location}, \mathrm{scale})$.

\begin{jwarn}
The \yamlkey{Log} distribution requires \yamlkey{min}$>0$: you cannot take the logarithm of zero
or a negative number. If part of your range must include zero or negatives, use \yamlkey{Flat}.
\end{jwarn}

\section{A multi-variable example}
\label{sec:var-example}

A scan usually mixes distributions---a flat mass here, a log-scaled coupling there:

\begin{yamlwin}
Sampling:
  Variables:
    - name: m0
      description: "Scalar mass (flat)"
      distribution: {type: Flat, parameters: {min: 100, max: 3000}}
    - name: tanb
      description: "tan(beta), log-scaled"
      distribution: {type: Log, parameters: {min: 1, max: 60}}
    - name: A0
      description: "Trilinear term (normal)"
      distribution: {type: Normal, parameters: {mean: 0, stddev: 1000}}
\end{yamlwin}

\begin{jnote}
All five distributions support \Jarvis's internal mapping from a uniform number in $[0,1]$, which
is what lets nested samplers and \textsc{mcmc} methods use them. Other families (Binomial,
Poisson, Beta, and so on) are not available as scan variables in the current version.
\end{jnote}

\section{Summary}
\label{sec:var-summary}

Each variable has a \yamlkey{name} and a \yamlkey{distribution} (a \yamlkey{type} plus its
\yamlkey{parameters}). Pick \yamlkey{Flat} for plain ranges, \yamlkey{Log} for wide positive
ranges, and the \yamlkey{Normal}/\yamlkey{Log-Normal}/\yamlkey{Logit} family when values should
concentrate around a centre. The variable names you choose here are exactly the symbols you will
use in the likelihood (Chapter~\ref{ch:loglikelihood}) and in expressions
(Chapter~\ref{ch:symbols}).
